\section{问题二的建模与求解}

本节围绕问题二的比较或检验目标，先统一比较口径，再建立主检验模型，最后结合稳健性证据给出结论。

\subsection{问题分析}
说明比较对象、处理因素、基准期、观测单位和需要避免的混淆。

\subsection{比较口径与数据构造}
说明面板、同店、基准期、变化量或其他比较指标的定义。

\subsection{基线检验与主检验模型}
\subsubsection{基线检验}
说明基础统计检验、效应量和区间。

\subsubsection{主检验模型}
给出模型设定、参数含义、标准误和显著性判断。

\subsection{稳健性检验与结果分析}
说明替代口径、窗口、分布或重采样检验，并解释结论是否一致。

\subsection{问题二小结}
用一句规模结论和一句变化/显著性结论直接回答问题二。
